\documentclass[12pt,a4paper]{article}

\usepackage[utf8]{inputenc}
\usepackage[T2A]{fontenc}
\usepackage[russian]{babel}
\usepackage{geometry}
\usepackage{setspace}
\usepackage{indentfirst}
\usepackage{amsmath}
\usepackage{amssymb}
\usepackage{array}
\usepackage{booktabs}
\usepackage{longtable}
\usepackage{float}
\usepackage{graphicx}
\usepackage{hyperref}
\usepackage{enumitem}
\usepackage{tikz}
\usetikzlibrary{arrows.meta,positioning,shapes.geometric,fit,calc}

\geometry{left=2.5cm,right=1.5cm,top=2cm,bottom=2cm}
\setstretch{1.2}
\setlength{\parindent}{1.25cm}
\setlist[itemize]{leftmargin=1.8cm}
\setlist[enumerate]{leftmargin=1.8cm}
\hypersetup{hidelinks}
\sloppy

\newcolumntype{L}[1]{>{\raggedright\arraybackslash}p{#1}}
\newcolumntype{C}[1]{>{\centering\arraybackslash}p{#1}}

\begin{document}

\begin{titlepage}
    \begin{center}
        \textbf{Министерство науки и высшего образования Российской Федерации}\\
        \textbf{Федеральное государственное бюджетное образовательное учреждение высшего образования}\\
        \textbf{<<\hspace{9cm}>>}\\[1cm]
        \textbf{Кафедра \hspace{8cm}}\\[3cm]

        {\large \textbf{КУРСОВАЯ РАБОТА}}\\[0.5cm]
        по дисциплине <<Архитектура вычислительных систем>>\\[0.8cm]
        на тему:\\[0.3cm]
        {\Large \textbf{Проектирование собственного TPU-подобного сопроцессора}}\\
        {\Large \textbf{для матричных и векторных операций}}\\[2.5cm]
    \end{center}

    \begin{flushright}
        Выполнил: студент группы \underline{\hspace{4cm}}\\[0.3cm]
        \underline{\hspace{5.2cm}}\\[0.8cm]
        Проверил: \underline{\hspace{5cm}}\\[0.3cm]
        \underline{\hspace{5.2cm}}
    \end{flushright}

    \vfill
    \begin{center}
        2026
    \end{center}
\end{titlepage}

\tableofcontents
\newpage

\section*{Введение}
\addcontentsline{toc}{section}{Введение}

Рост вычислительной сложности нейросетевых и сигнальных алгоритмов привёл к широкому распространению специализированных ускорителей, ориентированных на регулярные матричные и векторные вычисления. В отличие от универсального центрального процессора, такой ускоритель проектируется под узкий класс операций и поэтому обеспечивает более высокую эффективность по пропускной способности, энергопотреблению и детерминированности исполнения. Одним из наиболее известных примеров является Google TPU v1, в котором центральную роль играет крупный систолический массив, связанный с программно управляемой локальной памятью.

В рамках данной курсовой работы рассматривается проектирование собственного TPU-подобного сопроцессора как \textit{внешнего специализированного блока}, взаимодействующего с неизменяемым хост-процессором. Архитектура основного процессора не модифицируется: сопроцессор получает команды, данные и возвращает результаты через стандартный интерфейс обмена.

\textbf{Цель работы} --- разработать минималистичную, но детально описанную архитектуру сопроцессора, предназначенного исключительно для матричных и векторных операций: умножения матриц, свёрток, активаций и поэлементных преобразований.

\textbf{Для достижения цели поставлены следующие задачи:}
\begin{enumerate}
    \item проанализировать архитектурные принципы Google TPU v1 как прототипа;
    \item выбрать состав функциональных блоков собственного сопроцессора;
    \item разработать микроархитектуру вычислительного ядра, локальной памяти и управления;
    \item определить систему команд, формат дескрипторов и протокол обмена с хостом;
    \item описать прохождение типовых операций через все узлы сопроцессора;
    \item подготовить схемы архитектуры и взаимодействия в формате TikZ для последующей компиляции в \LaTeX.
\end{enumerate}

\textbf{Объект исследования} --- специализированный сопроцессор для тензорных вычислений.\\
\textbf{Предмет исследования} --- микроархитектура TPU-подобного ускорителя, его локальная память, командный интерфейс и алгоритм исполнения операций матричного типа.

\section{Глава 1. Анализ архитектуры Google TPU v1}

\subsection{Ключевые особенности Google TPU v1}

Google TPU v1 был разработан как специализированный ускоритель инференса нейронных сетей. Главная идея его архитектуры заключается в переносе вычислительно интенсивных операций умножения матриц в специализированный систолический массив, а также в использовании программно управляемой локальной памяти вместо традиционной многоуровневой кэш-иерархии.

С точки зрения архитектуры TPU v1 можно выделить следующие базовые принципы:
\begin{itemize}
    \item использование крупного систолического массива \(256 \times 256\), содержащего 65536 умножителей 8-битных чисел;
    \item накопление промежуточных результатов в более широком формате, что уменьшает потери точности;
    \item программно управляемый буфер данных (Unified Buffer), в который заранее подгружаются активации и промежуточные тензоры;
    \item потоковая подача весов и данных в вычислительное ядро;
    \item детерминированная модель исполнения без сложной спекуляции, переупорядочивания инструкций и аппаратной поддержки общего назначения.
\end{itemize}

Такое построение хорошо подходит для плотных линейных операций, поскольку умножение матриц и свёртки сводятся к регулярному шаблону <<умножение--накопление>>. В результате достигается высокая плотность вычислений на единицу площади кристалла при сравнительно простом управлении.

\subsection{Причины выбора TPU v1 в качестве основы}

Для учебного проекта архитектура TPU v1 удобна по нескольким причинам.

Во-первых, она хорошо декомпозируется на небольшое число блоков: интерфейс с хостом, локальная память, DMA, управляющее устройство, систолический массив и узел постобработки. Во-вторых, базовый поток данных легко описывается и визуализируется: данные загружаются в локальные буферы, затем последовательно проходят через массив, а результаты возвращаются в память хоста. В-третьих, такая архитектура масштабируема: уменьшение размера массива и буферов позволяет получить учебный вариант без разрушения исходной идеи.

Для курсовой работы особенно важно, что TPU-подобная организация позволяет:
\begin{itemize}
    \item подробно описать микроархитектуру отдельной вычислительной ячейки;
    \item формально задать систему команд, не зависящую от ISA хоста;
    \item выделить явный протокол обмена хост--сопроцессор;
    \item обосновать управление памятью через локальные банки и DMA-передачи.
\end{itemize}

\subsection{Ограничения оригинальной архитектуры и необходимость адаптации}

Несмотря на достоинства TPU v1, прямое копирование его параметров для учебного проекта нецелесообразно. Полноразмерный массив \(256 \times 256\) требует большого объёма локальной памяти, широкой внутренней шины и сложной системы питания и тактирования. Аналогично, инфраструктура серверного масштаба, ориентированная на реальные дата-центры, выходит за рамки курсовой работы.

Поэтому в работе предлагается \textbf{учебная микроархитектура TPU-Lite16}, сохраняющая принципы TPU v1, но имеющая два осознанных изменения:
\begin{enumerate}
    \item \textbf{уменьшенный систолический массив \(16 \times 16\)} вместо \(256 \times 256\), что радикально упрощает локальную память, маршрутизацию данных и конечный автомат управления;
    \item \textbf{добавленный генератор скользящего окна (Window Generator)} для выполнения \texttt{CONV2D} без полного развёртывания входа в \texttt{im2col} в памяти хоста.
\end{enumerate}

Кроме того, блок постобработки преднамеренно сделан облегченным: вместо отдельной сложной подсистемы нормализации и пуллинга предусмотрен компактный \textit{post-processing unit}, поддерживающий масштабирование, ограничение диапазона, ReLU/ReLU6 и LUT-аппроксимацию сигмоиды. Такое решение соответствует требованию минимализма и не выводит сопроцессор за пределы матрично-векторной специализации.

\begin{table}[H]
    \centering
    \caption{Сравнение TPU v1 и проектируемого сопроцессора}
    \begin{tabular}{|L{4.2cm}|L{4.7cm}|L{5.0cm}|}
        \hline
        \textbf{Параметр} & \textbf{Google TPU v1} & \textbf{TPU-Lite16} \\
        \hline
        Вычислительное ядро & Систолический массив \(256 \times 256\) & Систолический массив \(16 \times 16\) \\
        \hline
        Числовой формат & INT8 с накоплением в расширенном формате & INT8 с накоплением в INT32 \\
        \hline
        Локальная память & Крупная on-chip память серверного класса & Небольшие программно управляемые банки \\
        \hline
        Поддержка свёрток & Через преобразование в матричный вид & Через матричный режим и Window Generator \\
        \hline
        Подключение к хосту & Серверная платформа Google & Внешний PCIe-сопроцессор с DMA и MMIO \\
        \hline
        Постобработка & Специализированные блоки под инференс & Упрощённый блок активации и масштабирования \\
        \hline
    \end{tabular}
\end{table}

\section{Глава 2. Проектирование микроархитектуры сопроцессора}

\subsection{Общий замысел и основные параметры}

Проектируемый сопроцессор TPU-Lite16 предназначен для выполнения следующих операций:
\begin{itemize}
    \item блочное умножение матриц \texttt{MATMUL};
    \item двумерная свёртка \texttt{CONV2D};
    \item активации \texttt{ACTIVATE};
    \item поэлементные операции \texttt{EWISE} (сложение, умножение, максимум);
    \item загрузка и выгрузка тайлов между памятью хоста и локальными буферами.
\end{itemize}

Сопроцессор рассматривается как отдельное устройство на шине PCIe. Хост выделяет область памяти под дескрипторы и данные, а сопроцессор, используя DMA, самостоятельно переносит блоки данных в свою локальную память.

\begin{table}[H]
    \centering
    \caption{Базовые параметры микроархитектуры TPU-Lite16}
    \begin{tabular}{|L{5.0cm}|L{3.5cm}|L{6.0cm}|}
        \hline
        \textbf{Параметр} & \textbf{Значение} & \textbf{Назначение} \\
        \hline
        Размер систолического массива & \(16 \times 16\) PE & 256 операций MAC за такт в установившемся режиме \\
        \hline
        Формат входных операндов & INT8 & Минимизация аппаратной сложности \\
        \hline
        Формат аккумуляторов & INT32 & Безопасное накопление частичных сумм \\
        \hline
        Unified Buffer & 256 КиБ & Тайлы активаций, входов и промежуточных результатов \\
        \hline
        Weight Buffer & 64 КиБ & Хранение текущих весовых тайлов \\
        \hline
        Accumulator Buffer & 128 КиБ & Частичные суммы до постобработки \\
        \hline
        Таблица LUT & 4 КиБ & Аппроксимация сигмоиды/гиперболического тангенса \\
        \hline
        Интерфейс с хостом & PCIe + MMIO + DMA & Командный и данные-пути \\
        \hline
        Очередь команд & Кольцо дескрипторов по 64 байта & Передача задач от хоста \\
        \hline
    \end{tabular}
\end{table}

\subsection{Структурная схема сопроцессора}

На рисунке~\ref{fig:global-architecture} приведена общая структурная схема сопроцессора. Входным источником команд выступает хост-процессор, который размещает дескрипторы в памяти хоста. Блок PCIe/MMIO обеспечивает доступ к регистрам управления, а DMA Engine выполняет реальный перенос блоков данных. Управляющий автомат координирует предзагрузку буферов, запуск вычислительного ядра и запись результатов.

\begin{figure}[H]
    \centering
    \resizebox{\textwidth}{!}{
    \begin{tikzpicture}[
        block/.style={draw, rounded corners, minimum width=2.4cm, minimum height=1.1cm, align=center, fill=blue!8},
        mem/.style={draw, minimum width=2.7cm, minimum height=1.1cm, align=center, fill=green!10},
        ctrl/.style={draw, rounded corners, minimum width=2.7cm, minimum height=1.1cm, align=center, fill=orange!12},
        line/.style={-Latex, thick},
        dashedbox/.style={draw, dashed, rounded corners, inner sep=0.25cm}
    ]
        \node[block] (cpu) at (0,2.2) {Хост-процессор};
        \node[mem]   (dram) at (0,0) {Память хоста\\(дескрипторы, тензоры)};

        \node[block] (pcie) at (4.2,2.2) {PCIe Endpoint\\MMIO / Doorbell};
        \node[block] (dma)  at (4.2,0)   {DMA Engine\\Read / Write};

        \node[ctrl]  (queue) at (8.2,2.7) {Queue Manager\\Fetch / Decode};
        \node[ctrl]  (fsm)   at (8.2,1.1) {Control FSM\\Scoreboard};

        \node[mem]   (ub)    at (12.0,2.6) {Unified Buffer\\256 КиБ};
        \node[mem]   (wb)    at (12.0,1.1) {Weight Buffer\\64 КиБ};
        \node[mem]   (acc)   at (12.0,-0.4) {Accumulator Buffer\\128 КиБ};

        \node[block] (win)   at (15.8,2.2) {Window Generator\\для \texttt{CONV2D}};
        \node[block] (array) at (15.8,0.6) {Систолический массив\\\(16 \times 16\)};
        \node[block] (ppu)   at (19.4,0.6) {Post-Processing Unit\\Scale / ReLU / LUT};
        \node[block] (irq)   at (19.4,2.3) {IRQ / Status};

        \draw[line] (cpu) -- node[above] {MMIO} (pcie);
        \draw[line] (cpu) -- node[left] {запись дескрипторов} (dram);
        \draw[line] (pcie) -- (queue);
        \draw[line] (queue) -- (fsm);
        \draw[line] (dram) -- node[below] {DMA} (dma);
        \draw[line] (dma) -- (ub);
        \draw[line] (dma) -- (wb);
        \draw[line] (dma) -- (acc);
        \draw[line] (fsm) -- (ub);
        \draw[line] (fsm) -- (wb);
        \draw[line] (fsm) -- (acc);
        \draw[line] (ub) -- (win);
        \draw[line] (win) -- (array);
        \draw[line] (wb) -- (array);
        \draw[line] (array) -- (acc);
        \draw[line] (acc) -- (ppu);
        \draw[line] (ppu) -- (ub);
        \draw[line] (ppu) -- (dma);
        \draw[line] (fsm) -- (irq);
        \draw[line] (irq) -- node[above] {прерывание / статус} (cpu);

        \node[dashedbox, fit=(cpu)(dram)] {};
        \node[dashedbox, fit=(pcie)(dma)(queue)(fsm)(ub)(wb)(acc)(win)(array)(ppu)(irq)] {};

        \node at (0,-1.5) {\textbf{Сторона хоста}};
        \node at (12.2,-1.5) {\textbf{Сторона сопроцессора}};
    \end{tikzpicture}
    }
    \caption{Структурная схема TPU-Lite16}
    \label{fig:global-architecture}
\end{figure}

\textbf{Пояснение к рисунку.} Поток управления идёт сверху: хост через MMIO сообщает о новых командах, Queue Manager считывает дескрипторы и передаёт их в управляющий автомат. Поток данных идёт снизу: DMA переносит тайлы в локальную память, после чего данные проходят через Window Generator и систолический массив. Полученный результат поступает в буфер аккумуляторов и затем в блок постобработки.

\subsection{Систолический массив и вычислительная ячейка}

Базовым вычислительным узлом является PE (\textit{processing element}). Каждая ячейка принимает один элемент потока \(A\) слева и один элемент потока \(B\) сверху, умножает их и добавляет произведение к локальной накопленной сумме:
\[
S_{i,j}^{(t+1)} = S_{i,j}^{(t)} + A_{i,k} \cdot B_{k,j}.
\]

В проекте используется \textbf{output-stationary} схема: частичная сумма фиксируется в пределах PE до завершения обработки текущего тайла, а затем выгружается в \texttt{Accumulator Buffer}. Это упрощает маршрутизацию суммы и делает поведение массива детерминированным.

Одна PE содержит:
\begin{itemize}
    \item регистр входа \(A\) шириной 8 бит;
    \item регистр входа \(B\) шириной 8 бит;
    \item умножитель INT8\(\times\)INT8;
    \item сумматор с аккумулятором INT32;
    \item цепи передачи \(A\) вправо и \(B\) вниз;
    \item сигнал \texttt{clear\_acc} для обнуления перед новой задачей.
\end{itemize}

В установившемся режиме массив выполняет до 256 операций MAC за такт. Заполнение массива требует \(16 - 1 + 16 - 1 = 30\) тактов на распространение первой волны данных, после чего каждое последующее поступление входного элемента продвигает вычисление очередной диагонали.

\begin{figure}[H]
    \centering
    \resizebox{\textwidth}{!}{
    \begin{tikzpicture}[
        pe/.style={draw, minimum width=1.25cm, minimum height=0.9cm, fill=blue!8, align=center},
        cell/.style={draw, rounded corners, minimum width=0.95cm, minimum height=0.55cm, align=center, fill=orange!10},
        line/.style={-Latex, thick},
        note/.style={align=center}
    ]
        % Array fragment
        \node[pe] (p11) at (0,0) {PE(0,0)};
        \node[pe] (p12) at (1.8,0) {PE(0,1)};
        \node[pe] (p13) at (3.6,0) {PE(0,2)};
        \node[pe] (p21) at (0,-1.4) {PE(1,0)};
        \node[pe] (p22) at (1.8,-1.4) {PE(1,1)};
        \node[pe] (p23) at (3.6,-1.4) {PE(1,2)};
        \node[pe] (p31) at (0,-2.8) {PE(2,0)};
        \node[pe] (p32) at (1.8,-2.8) {PE(2,1)};
        \node[pe] (p33) at (3.6,-2.8) {PE(2,2)};

        \draw[line] (-1.3,0) -- node[above] {$A_0$} (p11);
        \draw[line] (-1.3,-1.4) -- node[above] {$A_1$} (p21);
        \draw[line] (-1.3,-2.8) -- node[above] {$A_2$} (p31);

        \draw[line] (0,1.0) -- node[left] {$B_0$} (p11);
        \draw[line] (1.8,1.0) -- node[left] {$B_1$} (p12);
        \draw[line] (3.6,1.0) -- node[left] {$B_2$} (p13);

        \foreach \from/\to in {p11/p12,p12/p13,p21/p22,p22/p23,p31/p32,p32/p33}
            \draw[line] (\from) -- (\to);
        \foreach \from/\to in {p11/p21,p21/p31,p12/p22,p22/p32,p13/p23,p23/p33}
            \draw[line] (\from) -- (\to);

        \node[note] at (1.8,-4.0) {Фрагмент массива \(16 \times 16\):\\поток \(A\) идёт слева направо, поток \(B\) --- сверху вниз};

        % PE internals
        \node[draw, rounded corners, minimum width=7.6cm, minimum height=4.0cm, fill=green!6] (box) at (11.0,-1.1) {};
        \node[note] at (11.0,0.9) {\textbf{Структура одной PE}};

        \node[cell] (areg) at (8.9,-0.2) {$A$ reg};
        \node[cell] (breg) at (8.9,-2.1) {$B$ reg};
        \node[cell] (mul)  at (10.8,-1.15) {$\times$};
        \node[cell] (add)  at (12.5,-1.15) {$+$};
        \node[cell] (accu) at (14.2,-1.15) {ACC\\INT32};

        \draw[line] (7.6,-0.2) -- node[above] {$A_{in}$} (areg);
        \draw[line] (7.6,-2.1) -- node[below] {$B_{in}$} (breg);
        \draw[line] (areg) -- (mul);
        \draw[line] (breg) -- (mul);
        \draw[line] (mul) -- (add);
        \draw[line] (add) -- (accu);
        \draw[line] (accu) to[out=90,in=0] node[above] {prev sum} (add);
        \draw[line] (areg) -- node[above] {$A_{out}$} (15.4,-0.2);
        \draw[line] (breg) -- node[below] {$B_{out}$} (15.4,-2.1);
        \draw[line] (accu) -- node[above] {$C_{tile}$} (15.4,-1.15);
    \end{tikzpicture}
    }
    \caption{Систолический массив и структура вычислительной ячейки}
    \label{fig:systolic-array}
\end{figure}

\textbf{Пояснение к рисунку.} Для наглядности показан только фрагмент массива \(3 \times 3\), но физически реализуется сетка \(16 \times 16\). Правая часть рисунка раскрывает внутреннее устройство PE: данные \(A\) и \(B\) регистрируются, участвуют в операции MAC и затем передаются соседям.

\subsection{Организация локальной памяти}

Локальная память программно управляется и не является кэшем. Это значит, что хост или драйвер должны явно задавать, какие тайлы загрузить и куда сохранить результаты. Такой подход выбран по двум причинам: во-первых, он соответствует философии TPU v1; во-вторых, он уменьшает сложность аппаратуры.

\begin{table}[H]
    \centering
    \caption{Организация локальных буферов}
    \begin{tabular}{|L{3.5cm}|L{2.3cm}|L{3.0cm}|L{6.0cm}|}
        \hline
        \textbf{Буфер} & \textbf{Объём} & \textbf{Банки} & \textbf{Назначение} \\
        \hline
        Unified Buffer & 256 КиБ & 16 банков по 16 КиБ & Входные матрицы, активации, временные тензоры, тайлы выходов \\
        \hline
        Weight Buffer & 64 КиБ & 8 банков по 8 КиБ & Хранение весов для \texttt{MATMUL} и \texttt{CONV2D} \\
        \hline
        Accumulator Buffer & 128 КиБ & 8 банков по 16 КиБ & Частичные суммы в INT32 до активации/квантизации \\
        \hline
        LUT Buffer & 4 КиБ & 1 банк & Таблицы аппроксимации нелинейностей \\
        \hline
    \end{tabular}
\end{table}

Адресация в банках выполняется по 16-байтным словам. Индекс банка вычисляется как
\[
\textit{bank} = \left(\frac{\textit{addr}}{16}\right) \bmod N_{banks},
\]
где \(N_{banks}\) --- число банков в конкретном буфере. Такое распределение уменьшает вероятность конфликтов при последовательном доступе.

Unified Buffer логически делится на две половины по 128 КиБ, что позволяет использовать \textit{ping-pong} режим: пока одна половина участвует в вычислении, во вторую DMA может подгружать следующий тайл. Аналогичный приём применим и к Weight Buffer.

Accumulator Buffer хранит результаты в формате INT32 и поддерживает два режима:
\begin{itemize}
    \item \textbf{accumulate mode} --- добавление нового частичного результата к уже существующему;
    \item \textbf{drain mode} --- передача блока в Post-Processing Unit с последующей записью квантованного или активированного результата.
\end{itemize}

\subsection{Подсистема активации и постобработки}

Чтобы не усложнять микроархитектуру, в сопроцессоре отсутствуют полнофункциональные блоки нормализации и пуллинга. Их роль выполняется частично программным способом, частично --- компактным блоком постобработки. Такой блок получает данные из \texttt{Accumulator Buffer} и выполняет одну из операций:
\begin{itemize}
    \item \texttt{identity};
    \item \texttt{ReLU};
    \item \texttt{ReLU6};
    \item \texttt{hard-sigmoid} по LUT или кусочно-линейной аппроксимации;
    \item \texttt{tanh} по LUT;
    \item насыщение и сдвиг при преобразовании INT32 \(\rightarrow\) INT8.
\end{itemize}

Batch Normalization в предлагаемой архитектуре не реализуется как отдельная инструкция, а \textbf{сворачивается в коэффициенты весов и смещений} на этапе подготовки модели. Это типичное упрощение для инференс-ускорителей.

\subsection{Генератор скользящего окна для свёрток}

Второе целевое изменение относительно упрощённого копирования TPU v1 --- это \textbf{Window Generator}. Он формирует последовательности адресов чтения из Unified Buffer таким образом, чтобы свёртка выполнялась без полного материализованного \texttt{im2col} в памяти хоста.

Блок получает параметры:
\begin{itemize}
    \item размер входного тензора \(H_{in}, W_{in}, C_{in}\);
    \item размер ядра \(K_h, K_w\);
    \item шаги \(stride_h, stride_w\);
    \item дополнение \(pad_t, pad_l, pad_b, pad_r\);
    \item дилатацию \(dilation_h, dilation_w\).
\end{itemize}

На каждом шаге генератор адресов выдаёт срез окна в формате, совместимом с подачей в левую границу систолического массива. За счёт этого свёртка математически сводится к матричному умножению, но физически лишний этап записи \texttt{im2col} в память отсутствует. Это снижает объём внешнего трафика и делает поддержку \texttt{CONV2D} более естественной.

\subsection{Интерфейс обмена с хостом}

Так как архитектура основного процессора по условию работы не меняется, сопроцессор проектируется как \textbf{периферийное PCIe-устройство}. Для обмена используются два канала:
\begin{enumerate}
    \item \textbf{MMIO-регистры} --- для настройки, указания адреса кольца дескрипторов, сигнала \texttt{doorbell}, чтения статуса и обработки прерываний;
    \item \textbf{DMA} --- для передачи тайлов матриц, векторов и результатов между памятью хоста и локальными буферами сопроцессора.
\end{enumerate}

\begin{table}[H]
    \centering
    \caption{Основные MMIO-регистры блока управления}
    \begin{tabular}{|C{2.0cm}|L{4.0cm}|L{7.2cm}|}
        \hline
        \textbf{Смещение} & \textbf{Регистр} & \textbf{Назначение} \\
        \hline
        0x00 & \texttt{VERSION} & Версия аппаратуры и поддерживаемый формат дескрипторов \\
        \hline
        0x04 & \texttt{CONTROL} & Сброс, запуск, маска прерываний, разрешение DMA \\
        \hline
        0x08 & \texttt{STATUS} & Состояние конечного автомата, код ошибки, признак занятости \\
        \hline
        0x10/0x14 & \texttt{QUEUE\_BASE\_LO/HI} & 64-битный адрес кольца дескрипторов в памяти хоста \\
        \hline
        0x18 & \texttt{QUEUE\_SIZE} & Размер кольца в дескрипторах \\
        \hline
        0x1C & \texttt{QUEUE\_HEAD} & Позиция последнего выполненного дескриптора \\
        \hline
        0x20 & \texttt{QUEUE\_TAIL} & Позиция, заполненная хостом \\
        \hline
        0x24 & \texttt{DOORBELL} & Сигнал о появлении новых команд \\
        \hline
        0x28 & \texttt{IRQ\_STATUS} & Источник прерывания: completion, fault, timeout \\
        \hline
        0x2C & \texttt{IRQ\_MASK} & Маскирование прерываний \\
        \hline
    \end{tabular}
\end{table}

Драйвер на стороне хоста обязан обеспечить выравнивание дескрипторов по 64 байта и размещение DMA-буферов в физически адресуемой памяти. Когерентность кэша не предполагается: перед постановкой задачи драйвер выполняет операции синхронизации буферов.

\subsection{Управляющее устройство и конвейер исполнения}

Управляющее устройство реализовано в виде конечного автомата с жёстко заданной последовательностью фаз:
\begin{enumerate}
    \item \texttt{IDLE} --- ожидание команды;
    \item \texttt{FETCH} --- считывание дескриптора из памяти хоста;
    \item \texttt{DECODE} --- декодирование полей, проверка диапазонов и зависимостей;
    \item \texttt{PREFETCH} --- запуск DMA или настройка генератора окна;
    \item \texttt{EXECUTE} --- запуск массива или блока постобработки;
    \item \texttt{WRITEBACK} --- запись результата в локальную память либо в память хоста;
    \item \texttt{COMPLETE} --- обновление \texttt{QUEUE\_HEAD}, установка статуса и генерация IRQ при необходимости.
\end{enumerate}

Дополнительно используется простой \textit{scoreboard}: он отслеживает, какие локальные области памяти заняты предыдущими командами, чтобы не допустить перезаписи тайла до завершения чтения. В отличие от процессоров общего назначения, переупорядочивание инструкций не поддерживается; команды исполняются \textbf{строго по порядку}, что упрощает верификацию.

\section{Глава 3. Система команд и интерфейс с хостом}

\subsection{Модель управления со стороны хоста}

Хост-процессор взаимодействует с TPU-Lite16 по следующему алгоритму:
\begin{enumerate}
    \item выделяет буферы в оперативной памяти под входные и выходные тензоры;
    \item формирует кольцо дескрипторов команд;
    \item записывает базовый адрес кольца в MMIO-регистры сопроцессора;
    \item размещает команды \texttt{LOAD}, \texttt{MATMUL}, \texttt{CONV2D}, \texttt{ACTIVATE}, \texttt{STORE}, \texttt{SYNC};
    \item обновляет хвост очереди \texttt{QUEUE\_TAIL} и подаёт \texttt{DOORBELL};
    \item ожидает завершения либо по прерыванию, либо опросом \texttt{STATUS}.
\end{enumerate}

Таким образом, основной процессор не управляет каждым тактом сопроцессора, а только задаёт \textit{грубозернистые операции}. Это соответствует модели автономного ускорителя.

\subsection{Формат дескриптора команды}

Все команды имеют единый размер \textbf{64 байта} (16 слов по 32 бита). Общий формат выбран намеренно: фиксированный размер дескриптора упрощает работу Queue Manager и уменьшает сложность декодера.

\begin{table}[H]
    \centering
    \caption{Заголовок дескриптора (\texttt{DW0})}
    \begin{tabular}{|C{2.7cm}|C{2.0cm}|L{8.0cm}|}
        \hline
        \textbf{Биты} & \textbf{Поле} & \textbf{Смысл} \\
        \hline
        7:0 & \texttt{OPCODE} & Код операции \\
        \hline
        11:8 & \texttt{DTYPE} & Формат данных: 0 --- INT8, 1 --- INT32 \\
        \hline
        15:12 & \texttt{FLAGS} & Флаги режима выполнения команды \\
        \hline
        23:16 & \texttt{CMD\_ID} & Идентификатор команды для трассировки и отладки \\
        \hline
        31:24 & \texttt{IRQ/FENCE} & Признак генерации IRQ и барьерные свойства \\
        \hline
    \end{tabular}
\end{table}

\small
\begin{longtable}{|C{1.2cm}|L{3.2cm}|L{9.2cm}|}
    \caption{Общий формат дескриптора команды} \\
    \hline
    \textbf{Слово} & \textbf{Поле} & \textbf{Назначение} \\
    \hline
    \endfirsthead
    \hline
    \textbf{Слово} & \textbf{Поле} & \textbf{Назначение} \\
    \hline
    \endhead
    DW0 & Заголовок & \texttt{OPCODE}, формат данных, флаги, ID команды \\
    \hline
    DW1 & \texttt{SRC0\_ADDR\_LO} & Младшие 32 бита первого адреса-источника \\
    \hline
    DW2 & \texttt{SRC0\_ADDR\_HI} & Старшие 32 бита первого адреса-источника \\
    \hline
    DW3 & \texttt{SRC1\_ADDR\_LO} & Младшие 32 бита второго адреса-источника \\
    \hline
    DW4 & \texttt{SRC1\_ADDR\_HI} & Старшие 32 бита второго адреса-источника \\
    \hline
    DW5 & \texttt{DST\_ADDR\_LO} & Младшие 32 бита адреса назначения \\
    \hline
    DW6 & \texttt{DST\_ADDR\_HI} & Старшие 32 бита адреса назначения \\
    \hline
    DW7 & \texttt{PARAM0} & Операционно-зависимый параметр 0 \\
    \hline
    DW8 & \texttt{PARAM1} & Операционно-зависимый параметр 1 \\
    \hline
    DW9 & \texttt{PARAM2} & Операционно-зависимый параметр 2 \\
    \hline
    DW10 & \texttt{PARAM3} & Операционно-зависимый параметр 3 \\
    \hline
    DW11 & \texttt{PARAM4} & Операционно-зависимый параметр 4 \\
    \hline
    DW12 & \texttt{PARAM5} & Операционно-зависимый параметр 5 \\
    \hline
    DW13 & \texttt{IMM0} & Дополнительные режимы или непосредственные значения \\
    \hline
    DW14 & \texttt{IMM1} & Дополнительные режимы или масштабирование \\
    \hline
    DW15 & \texttt{RESERVED} & Зарезервировано для будущего расширения \\
    \hline
\end{longtable}
\normalsize

Адреса в командах \texttt{LOAD} и \texttt{STORE} трактуются как 64-битные физические адреса памяти хоста. Адреса в вычислительных командах трактуются как \textbf{локальные адреса} внутри адресного пространства сопроцессора. Это адресное пространство разбивается следующим образом:
\begin{itemize}
    \item \texttt{0x0000\_0000 -- 0x0003\_FFFF} --- Unified Buffer;
    \item \texttt{0x0004\_0000 -- 0x0004\_FFFF} --- Weight Buffer;
    \item \texttt{0x0005\_0000 -- 0x0006\_FFFF} --- Accumulator Buffer;
    \item \texttt{0x0007\_0000 -- 0x0007\_0FFF} --- LUT Buffer и shadow-регистры.
\end{itemize}

\subsection{Набор команд}

\begin{table}[H]
    \centering
    \caption{Набор команд TPU-Lite16}
    \begin{tabular}{|C{1.8cm}|L{3.4cm}|L{8.3cm}|}
        \hline
        \textbf{Opcode} & \textbf{Команда} & \textbf{Назначение} \\
        \hline
        0x01 & \texttt{LOAD\_TILE} & Перенос тайла из памяти хоста в локальный буфер \\
        \hline
        0x02 & \texttt{STORE\_TILE} & Перенос тайла из локального буфера в память хоста \\
        \hline
        0x10 & \texttt{MATMUL} & Блочное умножение матриц на систолическом массиве \\
        \hline
        0x11 & \texttt{CONV2D} & Двумерная свёртка через Window Generator и массив \\
        \hline
        0x12 & \texttt{ACTIVATE} & Активация, масштабирование, насыщение \\
        \hline
        0x13 & \texttt{EWISE} & Поэлементные операции add/mul/max \\
        \hline
        0x14 & \texttt{SYNC} & Барьер, ожидание завершения всех предыдущих операций \\
        \hline
    \end{tabular}
\end{table}

\subsection{Формат команды \texttt{LOAD\_TILE}}

\begin{table}[H]
    \centering
    \caption{Поля команды \texttt{LOAD\_TILE}}
    \begin{tabular}{|L{2.6cm}|L{10.1cm}|}
        \hline
        \textbf{Поле} & \textbf{Содержимое} \\
        \hline
        \texttt{SRC0\_ADDR} & Физический адрес источника в памяти хоста \\
        \hline
        \texttt{DST\_ADDR} & Локальный адрес назначения (UB/WB/ACC) \\
        \hline
        \texttt{PARAM0} & \([15:0]\) --- число строк, \([31:16]\) --- число столбцов \\
        \hline
        \texttt{PARAM1} & \([15:0]\) --- шаг по строке в памяти хоста, \([31:16]\) --- шаг по строке в локальной памяти \\
        \hline
        \texttt{PARAM2} & \([7:0]\) --- размер элемента в байтах, \([15:8]\) --- длина DMA burst, \([23:16]\) --- код целевого буфера \\
        \hline
        \texttt{FLAGS} & bit0 --- генерация IRQ; bit1 --- дополнение нулями неполного хвоста; bit2 --- транспонирование при загрузке \\
        \hline
    \end{tabular}
\end{table}

Команда используется для подготовки тайлов матриц \(A\), \(B\), входных карт признаков и весов. При загрузке в Weight Buffer обычно запрещается транспонирование, а при загрузке в Unified Buffer оно может использоваться для удобного разложения матрицы по банкам.

\subsection{Формат команды \texttt{STORE\_TILE}}

\begin{table}[H]
    \centering
    \caption{Поля команды \texttt{STORE\_TILE}}
    \begin{tabular}{|L{2.6cm}|L{10.1cm}|}
        \hline
        \textbf{Поле} & \textbf{Содержимое} \\
        \hline
        \texttt{SRC0\_ADDR} & Локальный адрес источника (обычно UB или ACC после постобработки) \\
        \hline
        \texttt{DST\_ADDR} & Физический адрес результата в памяти хоста \\
        \hline
        \texttt{PARAM0} & \([15:0]\) --- число строк, \([31:16]\) --- число столбцов \\
        \hline
        \texttt{PARAM1} & \([15:0]\) --- локальный шаг по строке, \([31:16]\) --- шаг по строке в памяти хоста \\
        \hline
        \texttt{PARAM2} & \([7:0]\) --- размер элемента, \([15:8]\) --- burst, \([23:16]\) --- код исходного буфера \\
        \hline
        \texttt{FLAGS} & bit0 --- генерация IRQ; bit1 --- запись с насыщением; bit2 --- инвалидация локального буфера после выгрузки \\
        \hline
    \end{tabular}
\end{table}

\subsection{Формат команды \texttt{MATMUL}}

\begin{table}[H]
    \centering
    \caption{Поля команды \texttt{MATMUL}}
    \begin{tabular}{|L{2.6cm}|L{10.1cm}|}
        \hline
        \textbf{Поле} & \textbf{Содержимое} \\
        \hline
        \texttt{SRC0\_ADDR} & Базовый локальный адрес тайла матрицы \(A\) в Unified Buffer \\
        \hline
        \texttt{SRC1\_ADDR} & Базовый локальный адрес тайла матрицы \(B\) в Weight Buffer или Unified Buffer \\
        \hline
        \texttt{DST\_ADDR} & Базовый локальный адрес выходного тайла в Accumulator Buffer \\
        \hline
        \texttt{PARAM0} & \([7:0]\) --- активные строки \(M\), \([15:8]\) --- активные столбцы \(N\), \([31:16]\) --- общая длина свёртки \(K\) \\
        \hline
        \texttt{PARAM1} & \([15:0]\) --- шаг по строке матрицы \(A\) в байтах, \([31:16]\) --- шаг по строке матрицы \(B\) \\
        \hline
        \texttt{PARAM2} & \([15:0]\) --- шаг по строке выхода \(C\), \([31:16]\) --- число итераций по \(K\)-тайлам \\
        \hline
        \texttt{IMM0} & \([7:0]\) --- правый сдвиг результата, \([15:8]\) --- режим округления, \([31:16]\) --- зарезервировано \\
        \hline
        \texttt{FLAGS} & bit0 --- накопление к существующему \texttt{ACC}; bit1 --- очистка аккумулятора; bit2 --- \(A^T\); bit3 --- \(B^T\) \\
        \hline
    \end{tabular}
\end{table}

Команда вычисляет блочное произведение
\[
C_{M \times N} \leftarrow A_{M \times K} \cdot B_{K \times N}
\]
в пределах одного тайла. Для матриц большего размера хост формирует несколько команд \texttt{MATMUL} и управляет разбиением на блоки.

\subsection{Формат команды \texttt{CONV2D}}

\begin{table}[H]
    \centering
    \caption{Поля команды \texttt{CONV2D}}
    \begin{tabular}{|L{2.6cm}|L{10.1cm}|}
        \hline
        \textbf{Поле} & \textbf{Содержимое} \\
        \hline
        \texttt{SRC0\_ADDR} & Локальный адрес входной карты признаков в Unified Buffer \\
        \hline
        \texttt{SRC1\_ADDR} & Локальный адрес весов свёртки в Weight Buffer \\
        \hline
        \texttt{DST\_ADDR} & Локальный адрес выходного тайла в Accumulator Buffer \\
        \hline
        \texttt{PARAM0} & \([15:0]\) --- \(H_{in}\), \([31:16]\) --- \(W_{in}\) \\
        \hline
        \texttt{PARAM1} & \([15:0]\) --- \(C_{in}\), \([31:16]\) --- \(C_{out, tile}\) \\
        \hline
        \texttt{PARAM2} & \([7:0]\) --- \(K_h\), \([15:8]\) --- \(K_w\), \([23:16]\) --- \(stride_h\), \([31:24]\) --- \(stride_w\) \\
        \hline
        \texttt{PARAM3} & \([7:0]\) --- \(pad_t\), \([15:8]\) --- \(pad_l\), \([23:16]\) --- \(pad_b\), \([31:24]\) --- \(pad_r\) \\
        \hline
        \texttt{PARAM4} & \([15:0]\) --- \(H_{out}\), \([31:16]\) --- \(W_{out}\) \\
        \hline
        \texttt{PARAM5} & \([7:0]\) --- \(dilation_h\), \([15:8]\) --- \(dilation_w\), \([31:16]\) --- groups \\
        \hline
        \texttt{FLAGS} & bit0 --- накопление в текущий \texttt{ACC}; bit1 --- включить Window Generator; bit2 --- глубинная свёртка \\
        \hline
    \end{tabular}
\end{table}

Команда реализует вычисление
\[
Y[h,w,c_o] = \sum_{k_h=0}^{K_h-1} \sum_{k_w=0}^{K_w-1} \sum_{c_i=0}^{C_{in}-1}
X[h',w',c_i] \cdot W[k_h,k_w,c_i,c_o],
\]
где индексы \(h'\) и \(w'\) вычисляются с учётом шага, дополнения и дилатации. Window Generator подаёт нужные элементы \(X\) в таком порядке, чтобы вычислительное ядро использовало тот же путь, что и при \texttt{MATMUL}.

\subsection{Формат команды \texttt{ACTIVATE}}

\begin{table}[H]
    \centering
    \caption{Поля команды \texttt{ACTIVATE}}
    \begin{tabular}{|L{2.6cm}|L{10.1cm}|}
        \hline
        \textbf{Поле} & \textbf{Содержимое} \\
        \hline
        \texttt{SRC0\_ADDR} & Локальный адрес входного блока (обычно в Accumulator Buffer) \\
        \hline
        \texttt{DST\_ADDR} & Локальный адрес результата в Unified Buffer или обратно в Accumulator Buffer \\
        \hline
        \texttt{PARAM0} & \([15:0]\) --- строки, \([31:16]\) --- столбцы \\
        \hline
        \texttt{PARAM1} & \([15:0]\) --- входной шаг, \([31:16]\) --- выходной шаг \\
        \hline
        \texttt{IMM0} & \([7:0]\) --- код активации: 0=identity, 1=ReLU, 2=ReLU6, 3=hard-sigmoid, 4=tanh \\
        \hline
        \texttt{IMM1} & \([15:0]\) --- коэффициент масштаба, \([31:16]\) --- правый сдвиг/квантование \\
        \hline
        \texttt{FLAGS} & bit0 --- насыщение в INT8; bit1 --- использовать LUT; bit2 --- запись INT32 без квантизации \\
        \hline
    \end{tabular}
\end{table}

\subsection{Формат команды \texttt{EWISE}}

\begin{table}[H]
    \centering
    \caption{Поля команды \texttt{EWISE}}
    \begin{tabular}{|L{2.6cm}|L{10.1cm}|}
        \hline
        \textbf{Поле} & \textbf{Содержимое} \\
        \hline
        \texttt{SRC0\_ADDR} & Локальный адрес первого операнда \\
        \hline
        \texttt{SRC1\_ADDR} & Локальный адрес второго операнда либо не используется при скалярном режиме \\
        \hline
        \texttt{DST\_ADDR} & Локальный адрес результата \\
        \hline
        \texttt{PARAM0} & \([15:0]\) --- строки, \([31:16]\) --- столбцы \\
        \hline
        \texttt{PARAM1} & \([15:0]\) --- шаг \texttt{SRC0}, \([31:16]\) --- шаг \texttt{SRC1} \\
        \hline
        \texttt{PARAM2} & \([15:0]\) --- шаг \texttt{DST}, \([23:16]\) --- код операции (0=add, 1=mul, 2=max) \\
        \hline
        \texttt{IMM0} & Скалярное значение при флаге \texttt{scalar\_mode} \\
        \hline
        \texttt{FLAGS} & bit0 --- второй операнд скаляр; bit1 --- насыщение; bit2 --- беззнаковый режим \\
        \hline
    \end{tabular}
\end{table}

\subsection{Формат команды \texttt{SYNC}}

\begin{table}[H]
    \centering
    \caption{Поля команды \texttt{SYNC}}
    \begin{tabular}{|L{2.6cm}|L{10.1cm}|}
        \hline
        \textbf{Поле} & \textbf{Содержимое} \\
        \hline
        \texttt{IMM0} & Маска ожидаемых блоков: DMA, массив, PPU, очередь команд \\
        \hline
        \texttt{IMM1} & Тайм-аут ожидания в тактах или микросекундах \\
        \hline
        \texttt{FLAGS} & bit0 --- сгенерировать IRQ по завершении; bit1 --- очистить код ошибки \\
        \hline
    \end{tabular}
\end{table}

Команда \texttt{SYNC} не выполняет вычислений, а гарантирует, что все ранее выданные операции завершены и их результаты видимы последующим командам или хосту.

\subsection{Типовая последовательность взаимодействия хост--сопроцессор}

Для операции умножения матриц \(C = A \times B\) взаимодействие хоста и сопроцессора имеет следующий вид:
\begin{enumerate}
    \item хост подготавливает тайлы матриц \(A\) и \(B\) в памяти;
    \item хост помещает в кольцо команд \texttt{LOAD\_TILE(A)}, \texttt{LOAD\_TILE(B)}, \texttt{MATMUL}, \texttt{ACTIVATE} (при необходимости), \texttt{STORE\_TILE(C)}, \texttt{SYNC};
    \item Queue Manager считывает команды и запускает DMA;
    \item DMA переносит тайлы в Unified Buffer и Weight Buffer;
    \item массив вычисляет результат и пишет его в Accumulator Buffer;
    \item блок постобработки при необходимости выполняет активацию;
    \item результат выгружается в память хоста;
    \item сопроцессор выставляет признак завершения и формирует прерывание.
\end{enumerate}

\begin{figure}[H]
    \centering
    \resizebox{\textwidth}{!}{
    \begin{tikzpicture}[
        line/.style={-Latex, thick},
        lifeline/.style={gray, dashed},
        every node/.style={font=\small}
    ]
        \node at (0,0.5) {\textbf{Хост CPU}};
        \node at (4.5,0.5) {\textbf{Память хоста}};
        \node at (9.0,0.5) {\textbf{DMA / Queue}};
        \node at (13.5,0.5) {\textbf{Вычислительное ядро}};

        \draw[lifeline] (0,0) -- (0,-10.5);
        \draw[lifeline] (4.5,0) -- (4.5,-10.5);
        \draw[lifeline] (9.0,0) -- (9.0,-10.5);
        \draw[lifeline] (13.5,0) -- (13.5,-10.5);

        \draw[line] (0,-0.8) -- node[above] {1. запись дескрипторов и данных} (4.5,-0.8);
        \draw[line] (0,-1.8) -- node[above] {2. MMIO: \texttt{QUEUE\_TAIL}, \texttt{DOORBELL}} (9.0,-1.8);
        \draw[line] (9.0,-2.8) -- node[above] {3. fetch descriptor} (4.5,-2.8);
        \draw[line] (9.0,-3.8) -- node[above] {4. DMA read: тайлы \(A,B\)} (4.5,-3.8);
        \draw[line] (9.0,-4.8) -- node[above] {5. старт \texttt{MATMUL}} (13.5,-4.8);
        \draw[line] (13.5,-6.0) -- node[above] {6. результат в ACC / PPU} (9.0,-6.0);
        \draw[line] (9.0,-7.2) -- node[above] {7. DMA write: тайл \(C\)} (4.5,-7.2);
        \draw[line] (9.0,-8.4) -- node[above] {8. completion / status} (0,-8.4);
        \draw[line] (0,-9.4) -- node[above] {9. чтение результата} (4.5,-9.4);
    \end{tikzpicture}
    }
    \caption{Последовательность взаимодействия хост--сопроцессор}
    \label{fig:host-sequence}
\end{figure}

\textbf{Пояснение к рисунку.} На временной диаграмме видно, что после постановки команд хост практически не участвует в исполнении, пока сопроцессор не сообщит о завершении. Это разгружает центральный процессор и позволяет сопроцессору работать автономно.

\section{Глава 4. Реализация основных операций}

\subsection{Прохождение команды \texttt{MATMUL}}

Рассмотрим вычисление матричного произведения \(C = A \times B\), где
\[
A \in \mathbb{Z}_8^{M \times K}, \qquad B \in \mathbb{Z}_8^{K \times N}, \qquad C \in \mathbb{Z}_{32}^{M \times N}.
\]

Поскольку массив имеет размер \(16 \times 16\), большие матрицы разбиваются на тайлы размером не более \(16 \times 16\) по измерениям \(M\) и \(N\), а измерение \(K\) обрабатывается серией проходов.

\textbf{Последовательность исполнения:}
\begin{enumerate}
    \item Хост выдаёт \texttt{LOAD\_TILE} для блока \(A_{tile}\) в Unified Buffer.
    \item Хост выдаёт \texttt{LOAD\_TILE} для блока \(B_{tile}\) в Weight Buffer.
    \item \texttt{MATMUL} очищает либо переиспользует соответствующий блок в Accumulator Buffer.
    \item Данные \(A\) подаются в левую границу массива построчно, данные \(B\) --- в верхнюю границу по столбцам.
    \item Внутри массива выполняется волновое распространение и накопление частичных сумм.
    \item После завершения всех проходов по измерению \(K\) результат тайла выгружается в Accumulator Buffer.
    \item Команда \texttt{ACTIVATE} либо оставляет данные в INT32, либо преобразует их в INT8 с выбранной функцией активации.
    \item Команда \texttt{STORE\_TILE} переносит выходной тайл в память хоста.
\end{enumerate}

С точки зрения контроллера эта операция реализуется конвейерно: пока массив считает текущий тайл, DMA может подготавливать следующий блок в свободную половину Unified Buffer. Именно поэтому разбиение буферов на две половины является важным элементом архитектуры, даже при общей минималистичности системы.

\subsection{Прохождение команды \texttt{CONV2D}}

Свёртка
\[
Y = X * W
\]
в TPU-Lite16 исполняется тем же массивом, что и \texttt{MATMUL}. Различие заключается только в способе формирования входного потока \(X\). Вместо явного построения матрицы \texttt{im2col} Window Generator генерирует адреса элементов окна <<на лету>>.

\textbf{Последовательность исполнения \texttt{CONV2D}:}
\begin{enumerate}
    \item входной тензор \(X\) загружается в Unified Buffer;
    \item весовые коэффициенты \(W\) загружаются в Weight Buffer;
    \item декодер \texttt{CONV2D} программирует Window Generator параметрами ядра, шагом, паддингом и дилатацией;
    \item Window Generator поочерёдно формирует срезы входного тензора, соответствующие каждому положению свёрточного окна;
    \item сформированный поток поступает в левую границу массива, веса --- в верхнюю;
    \item частичные суммы пишутся в Accumulator Buffer;
    \item при необходимости выполняются \texttt{ACTIVATE} и последующая \texttt{STORE\_TILE}.
\end{enumerate}

Такой способ полезен по двум причинам. Во-первых, уменьшается объём обращений к памяти хоста: не нужно сохранять промежуточную развёртку. Во-вторых, один и тот же массив используется и для матричного умножения, и для свёртки, то есть вычислительное ядро не дублируется.

\subsection{Реализация активаций и поэлементных операций}

Команда \texttt{ACTIVATE} применяется после накопления результатов. Наиболее типичный сценарий --- преобразование INT32 результата в INT8:
\[
y = \text{clip}\left(\frac{x \cdot scale}{2^{shift}}\right),
\]
после чего к значению применяется выбранная нелинейность. В случае \texttt{ReLU} достаточно сравнения с нулём. Для \texttt{hard-sigmoid} и \texttt{tanh} используется таблица LUT, заранее загруженная драйвером.

Команда \texttt{EWISE} нужна для операций, которые не требуют массива:
\begin{itemize}
    \item добавление bias к вектору или матрице;
    \item поэлементное умножение на коэффициенты масштабирования;
    \item вычисление \texttt{max} для простейших вариантов редукции и квазипуллинга.
\end{itemize}

Таким образом, даже без выделенного блока нормализации сопроцессор сохраняет достаточную функциональность для инференс-задач малой и средней сложности.

\subsection{Пример командной последовательности для умножения матриц}

Пусть требуется вычислить
\[
C_{32 \times 32} = A_{32 \times 32} \times B_{32 \times 32}.
\]
Тогда хост разобьёт задачу на четыре выходных тайла \(16 \times 16\), для каждого из которых последовательно сформирует:
\begin{enumerate}
    \item \texttt{LOAD\_TILE} для соответствующего блока \(A\);
    \item \texttt{LOAD\_TILE} для соответствующего блока \(B\);
    \item \texttt{MATMUL} с \(M=16\), \(N=16\), \(K=16\);
    \item повторный \texttt{MATMUL} для второй половины измерения \(K\) с флагом накопления;
    \item \texttt{ACTIVATE} в режиме \texttt{identity} или \texttt{ReLU};
    \item \texttt{STORE\_TILE};
    \item \texttt{SYNC} в конце всей последовательности.
\end{enumerate}

Если вместо матрицы-матрицы требуется матрица-вектор, то используется тот же формат, но \(N=1\). Тем самым аппаратная специализация не усложняется дополнительными режимами: вектор рассматривается как частный случай матрицы.

\section*{Заключение}
\addcontentsline{toc}{section}{Заключение}

В данной курсовой работе спроектирован TPU-подобный сопроцессор TPU-Lite16, предназначенный для матричных и векторных операций и взаимодействующий с неизменяемым хост-процессором. За основу была взята архитектурная идея Google TPU v1: систолический массив, программно управляемая локальная память, детерминированное управление и вынесение сложных функций на уровень подготовки команд.

В результате проектирования определены:
\begin{itemize}
    \item состав аппаратных блоков сопроцессора;
    \item микроархитектура систолического массива \(16 \times 16\);
    \item объёмы и организация Unified Buffer, Weight Buffer и Accumulator Buffer;
    \item интерфейс обмена с хостом через PCIe, MMIO и DMA;
    \item формат 64-байтного командного дескриптора;
    \item набор команд \texttt{LOAD\_TILE}, \texttt{STORE\_TILE}, \texttt{MATMUL}, \texttt{CONV2D}, \texttt{ACTIVATE}, \texttt{EWISE}, \texttt{SYNC};
    \item последовательность исполнения операций матричного умножения и свёртки.
\end{itemize}

Ключевыми обоснованными отличиями от TPU v1 стали уменьшение массива до учебного масштаба и добавление генератора скользящего окна для прямой поддержки свёрток без полного \texttt{im2col}. Предложенная архитектура остаётся минималистичной, но при этом достаточно детализирована для использования в качестве законченного проекта курсовой работы и дальнейшего перехода к RTL-моделированию.

\section*{Список литературы}
\addcontentsline{toc}{section}{Список литературы}

\begin{enumerate}
    \item Jouppi N. P., Young C., Patil N. et al. In-Datacenter Performance Analysis of a Tensor Processing Unit // Proceedings of the 44th Annual International Symposium on Computer Architecture (ISCA). --- 2017.
    \item Patterson D., Hennessy J. Computer Organization and Design: The Hardware/Software Interface. --- 5th ed. --- Morgan Kaufmann, 2014.
    \item Sze V., Chen Y.-H., Yang T.-J., Emer J. Efficient Processing of Deep Neural Networks: A Tutorial and Survey // Proceedings of the IEEE. --- 2017. --- Vol. 105, No. 12.
    \item Zhang C., Li P., Sun G. et al. Optimizing FPGA-based Accelerator Design for Deep Convolutional Neural Networks // FPGA'15. --- 2015.
    \item PCI Express Base Specification [Электронный ресурс]. --- Режим доступа: \url{https://pcisig.com/specifications/pciexpress}.
    \item AXI Reference Guide / AMBA Specifications [Электронный ресурс]. --- Режим доступа: \url{https://developer.arm.com/architectures/system-architectures/amba/amba-specifications}.
    \item Классические учебные материалы по систолическим массивам, локальной памяти ускорителей и проектированию DMA-контроллеров --- источник для последующего дополнения студентом в соответствии с требованиями кафедры.
\end{enumerate}

\end{document}
